% --- Document: Preamble ---
\documentclass[../main.tex]{subfiles}
\graphicspath{{\subfix{../assets/}}}
% --- Document: Start ---
\begin{document}
%implementation
\section{Resultados Esperados}
apresentar quais resultados esperamos alcançar com o desenvolvimento do projeto.


% Automoção de atividades 
\subsection{Automação de atividades} A automação de atividades é essencial para a gestão eficiente de processos internos, pois reduz 
o tempo gasto com tarefas manuais, diminui a ocorrência de erros humanos e aumenta a confiabilidade das informações. Automatizar atividades 
repetitivas permite que a equipe se concentre em tarefas mais estratégicas, agregando mais valor ao negócio.

Segundo Gonçalves (2004), a automação de processos é uma ferramenta fundamental para assegurar a continuidade dos negócios e a 
confiabilidade das informações, minimizando os riscos de falhas causadas por erro humano ou esquecimento. A ABNT (2003) também 
destaca, através de suas normas de segurança da informação, a importância de manter a integridade e a precisão dos dados, princípios 
que a automação reforça.

    \subsubsection{Atividades repetitivas} 
    Atividades que são realizadas frequentemente, como geração de relatórios, contratos e recibos simples, serão automatizadas. Com 
    isso, garante-se uniformidade na emissão de documentos e ganho de eficiência operacional.

    \subsubsection{Atualizações obrigatórias} 
    O sistema controlará atualizações importantes, como mudanças de conta bancária, telefone e e-mail, além de reajustes de contratos. 
    Todo o histórico será mantido, permitindo rastrear alterações e suas respectivas motivações, garantindo transparência e segurança 
    nos registros.

    \subsubsection{Geração de termos e anexos} 
    A criação de termos contratuais e outros documentos poderá ser feita automaticamente, com poucos cliques. Além disso, será 
    possível anexar documentos de apoio, assegurando organização e facilidade de acesso às informações.

    \subsubsection{Evitar erros humanos} 
    A automação proporciona maior precisão no controle e manipulação dos dados, evitando erros na geração de recibos, atualizações 
    cadastrais, contratos, relatórios ou reajustes de valores. Isso contribui diretamente para a melhoria da qualidade das informações 
    e para a credibilidade do sistema.

    \subsubsection{Alertas e pendências} 
    O sistema gerará alertas e pendências internos para os usuários, informando sobre vencimentos de contas, encerramento de contratos, 
    necessidade de renovação de garantias e outros prazos importantes. Assim, evita-se o esquecimento de tarefas críticas e assegura-se 
    que todas as obrigações sejam cumpridas em tempo hábil.

    
    % Disponibilidade
\subsection{Disponibilidade}
A disponibilidade de um software refere-se à sua capacidade de estar operacional e acessível sempre que necessário. 
Segundo Bass, Clements e Kazman (2012), disponibilidade é definida como a habilidade de um sistema continuar operando 
mesmo na presença de falhas. É um dos atributos de qualidade mais críticos, pois o sistema precisará estar disponível
para os usuários a qualquer momento. Por se tratar de uma proposta de um servidor local, será implementada a medida 
de redundância do tipo RAID 1.
    \subsubsection{Rede interna e servidor local}
    Nossa proposta para garantir a disponibilidade do sistema consiste em implementá-lo dentro de uma rede local (LAN), 
    limitando o acesso apenas aos computadores conectados à mesma rede do servidor. Desta forma, o sistema não dependerá 
    da conexão com a internet para funcionar, o que elimina problemas relacionados a instabilidades externas e riscos 
    adicionais de segurança provenientes da exposição pública.
    \subsubsection{Redundância}
    Além disso, o servidor contará com um sistema de armazenamento redundante baseado em RAID (Redundant Array of Independent Disks), 
    utilizando três discos rígidos de 500 GB cada. Segundo Patterson, Gibson e Katz (1988), o RAID é uma técnica que combina múltiplos 
    discos físicos em uma única unidade lógica para melhorar a confiabilidade e o desempenho. No nosso caso, a utilização de RAID 
    proporcionará tolerância a falhas: caso um dos discos apresente defeito, o sistema continuará operando normalmente, evitando a 
    interrupção do serviço.

    A importância do RAID para a disponibilidade é amplamente reconhecida na arquitetura de software. Como apontado por Chen \textit{et al.} (1994), 
    técnicas de RAID são fundamentais para proteger dados contra falhas de hardware, reduzir tempos de recuperação e aumentar a 
    confiança no armazenamento de dados críticos.

\subsection{Backup}
O backup é um procedimento fundamental para garantir a segurança da informação, possibilitando a recuperação de dados em caso de falhas, 
perdas acidentais, ataques cibernéticos ou desastres físicos. Segundo Bass, Clements e Kazman (2012), manter cópias de segurança é uma 
prática essencial para assegurar a continuidade dos serviços de software. A ausência de um sistema de backup eficiente pode acarretar 
prejuízos irreparáveis para a organização, incluindo a perda de dados críticos e paralisação de operações.

Conforme orienta a Associação Brasileira de Normas Técnicas (ABNT, 2003), a proteção dos dados deve ser parte integrante da estratégia 
de manutenção da integridade da informação, sendo recomendada a adoção de múltiplas cópias armazenadas em diferentes locais.
    \subsubsection{Metodologia do Backup}

    A metodologia adotada prevê a realização de backups manuais diariamente, ao final do expediente. O processo será iniciado apenas 
    após a confirmação de que não haverá mais alterações nos dados do sistema naquele dia. Tal abordagem visa garantir que todas as 
    informações relevantes estejam devidamente atualizadas antes da criação das cópias de segurança.

    Segundo Gonçalves (2004), backups diários aumentam significativamente a proteção contra perdas de dados recentes, minimizando 
    riscos operacionais.

    \subsubsection{Backup Local}

    O backup local será realizado por meio da cópia completa do banco de dados para dispositivos físicos, especificamente um 
    pendrive e um disco rígido externo (HD). Essa prática visa assegurar uma recuperação rápida em caso de falha do sistema principal, 
    conforme orientações de Patterson, Gibson e Katz (1988), que destacam a importância de múltiplas camadas de proteção de dados para 
    garantir alta disponibilidade e confiabilidade.

    O armazenamento físico próximo ao local de operação permite respostas imediatas a eventuais incidentes, além de oferecer 
    independência de conexões externas.

    \subsubsection{Backup em Nuvem}

    Além do backup local, cópias do banco de dados serão enviadas para dois serviços distintos de armazenamento em nuvem, o 
    Google Drive e o OneDrive. Essa abordagem segue as boas práticas recomendadas por Chen et al. (1994), que destacam a 
    importância da redundância geográfica para proteger dados contra desastres locais.

    A utilização de múltiplos provedores em nuvem proporciona maior resiliência, pois em caso de falha ou indisponibilidade 
    em um dos serviços, as cópias de segurança ainda estarão acessíveis no segundo serviço.

    \subsubsection{Disponibilidade do Backup}

    Os backups serão armazenados por um período contínuo de 30 dias. Isso significa que a cada dia será criado um novo 
    backup, e o backup mais antigo, ao completar 30 dias, será excluído. Esse método, conhecido como retenção rotativa, 
    mantém uma janela de restauração ampla, permitindo a recuperação de dados de qualquer dia do último mês.

    De acordo com Bass, Clements e Kazman (2012), a definição de uma política clara de retenção de backups é fundamental 
    para equilibrar a disponibilidade de dados históricos e a gestão eficiente de armazenamento.
    
    \subsubsection{Referências para salvar em outro arquivo}
    Bass, L., Clements, P., \& Kazman, R. (2012). Software Architecture in Practice (3rd ed.). Addison-Wesley.

    Patterson, D. A., Gibson, G., \& Katz, R. H. (1988). A Case for Redundant Arrays of Inexpensive Disks (RAID). Proceedings of the 1988 ACM SIGMOD International Conference on Management of Data.

    Chen, P. M., Lee, E. K., Gibson, G. A., Katz, R. H., \& Patterson, D. A. (1994). RAID: High-Performance, Reliable Secondary Storage. ACM Computing Surveys (CSUR).
    
    ABNT, Associação Brasileira de Normas Técnicas. NBR ISO/IEC 17799: Código de prática para a gestão da segurança da informação. Rio de Janeiro: ABNT, 2003.

    Gonçalves, Vicente de Paula. Gerenciamento da continuidade dos negócios: uma abordagem prática. Rio de Janeiro: Qualitymark, 2004.
    
    DATE, C. J. Introdução a sistemas de banco de dados. 8. ed. Rio de Janeiro: Campus, 2004.

    HELMAN, P. NoSQL: Banco de Dados para Aplicações Web. São Paulo: Novatec, 2015.

    TURNBULL, J. The Docker Book: Containerization is the New Virtualization. 1. ed. [S.l.]: James Turnbull, 2014.
    FOWLER, M. Patterns of Enterprise Application Architecture. Boston: Addison-Wesley, 2002

    TILKOV, S.; VOGELS, W. Node.js: Using JavaScript to Build High-Performance Network Programs. IEEE Internet Computing, v. 14, n. 6, p. 80-83, 2010.

    GAMMA, E. et al. Design Patterns: Elements of Reusable Object-Oriented Software. Boston: Addison-Wesley, 1995.

    FIELDING, R. T. Architectural Styles and the Design of Network-based Software Architectures. 2000. Tese (Doutorado em Ciência da Computação) – University of California, Irvine.

    MICROSOFT. TypeScript Documentation. 2012. Disponível em: https://www.typescriptlang.org/. Acesso em: 28 abr. 2025.

    HOLMES, A. Getting MEAN with Mongo, Express, Angular, and Node. Shelter Island: Manning Publications, 2016.

    KARGER, A. Mongoose for Application Development. Birmingham: Packt Publishing, 2013.

    TAILWIND LABS. Tailwind CSS Documentation. 2017. Disponível em: https://tailwindcss.com/docs. Acesso em: 28 abr. 2025.
    
    NPM. bcrypt. Disponível em: https://www.npmjs.com/package/bcrypt. Acesso em: 28 abr. 2025.

    NPM. dotenv. Disponível em: https://www.npmjs.com/package/dotenv. Acesso em: 28 abr. 2025.

    NPM. jsonwebtoken. Disponível em: https://www.npmjs.com/package/jsonwebtoken. Acesso em: 28 abr. 2025.
    % --- Document: End ---
\end{document}

